% OWNER: sections/s2_2_auth.tex. Edit only this file.
\seclabel
\claim{\secnum{2.2}{sec:auth}\quad Access tokens last 900\,s; neither renewal nor logout changes the 7-day cookie jar}

\lead{Logging in returns the access token and the list of portfolios in one reply.} \code{POST /api/Auth/Login} with \code{{"email", "password"}} returns HTTP 200, an access token of 435 characters in JWT format (\code{accessToken}), 6 portfolios (\code{pfs}), 16 competitions (\code{ctsts}) and the server's current business day (\code{currBizDay}). The same reply carries personal fields (contact details), so the wrapper keeps only the fields it needs and never logs this reply in full.

\lead{The login reply sets 5 credentials with different lifetimes, and the wrapper needs only the access token and the cookie jar (\figref{creds}).} The access token goes in the \code{Authorization: Bearer} header of every REST request and of the feed's negotiation. The \code{refreshToken} cookie, sent together with the other root-path cookies, obtains a new access token.

% OWNER: figs/creds.tex. Edit only this file.
\fig[wide]{creds}{The access token lasts 900\,s; the wrapper stores the whole cookie jar}[Credentials returned by a login at 05:38:54 KST on 24 Nov 2026, and where the wrapper stores each]{%
{\tabfont\begin{tabularx}{\linewidth}{>{\RaggedRight\arraybackslash}p{24mm} >{\RaggedRight\arraybackslash}p{20mm} >{\RaggedRight\arraybackslash}p{50mm} >{\RaggedRight\arraybackslash}p{20mm} >{\RaggedRight\arraybackslash}X}
\thd{Name} & \thd{Carried in} & \thd{Lifetime} & \thd{Stored in} & \thd{Used for}\\
\hl\code{accessToken} & Response body & \textbf{900\,s}, JWT \code{exp} $-$ \code{iat}; claims include \code{Id}, \code{exp}, \code{iat}, \code{nbf}, \code{role} & Memory only & \code{Authorization: Bearer} header on every REST request\\\rowsep
\hl\code{refreshToken} & Cookie & \textbf{7 days}, to 1 Dec 05:38; path \code{/} & \code{session.json} & Renewal through \code{Auth/SilentRefresh}, sent with the other root-path cookies\\\rowsep
\code{asset_token} & Cookie & 1 hour; path \code{/assets/app/} & \code{session.json} & Downloading the page code served only to logged-in users\\\rowsep
\code{userid} & Cookie & 7 days & \code{session.json} & User ID read by the page's JavaScript\\\rowsep
\code{nxdid} & Cookie & 1 year & \code{session.json} & Probably a device identifier\\
\end{tabularx}}}{Shaded rows are the 2 credentials whose lifetimes set the wrapper's schedule: the first REST request more than 720\,s after the access token was issued (server clock) renews it, and the wrapper logs in again when \code{refreshToken} has less than 1 day left. Renewal with \code{refreshToken} alone was refused with a 401, so the wrapper writes every cookie of the jar to \code{~/.config/pytimefolio/session.json} (mode 0600) after login and after each renewal; the access token is never written to disk. Cookie values were not recorded, only names, paths and lifetimes. The purpose of \code{nxdid} is inferred from its name and lifetime.}{Server response headers; \code{pytimefolio/auth.py}; DART Quant Team}


\lead{Renewal needs the whole cookie jar, and it leaves the 7-day deadline where it was.} \code{POST /api/Auth/SilentRefresh} with the root-path cookies and no \code{Authorization} header returns HTTP 200 and a new access token. Sending \code{refreshToken} alone was refused with a 401, so the wrapper keeps and sends the whole cookie jar (D4). The expiry of \code{refreshToken} was 20:38:54 UTC on 30 Nov both before and after the call, and the reply set no new \code{refreshToken}. Renewal therefore neither slides the expiry forward nor rotates the token, and repeated renewals with the same cookie all succeeded. During the trading session 2 \code{SilentRefresh} calls sent at the same moment with the same cookies both returned 200, a third sent right after also returned 200, and none set a new \code{refreshToken}; concurrent renewals do not invalidate each other.

\lead{The server accepts an expired access token for 300\,s.} A first test bracketed the cut-off between 3.5 and 5.6\,min after expiry. A second token, issued at 08:16:43 by the server's clock, was sent again just before and just after 5 minutes past its expiry. By the server's \code{Date} header it was accepted 4~min~57\,s after expiry and refused 5~min~07\,s after (\figref{token}). The cut-off lies inside that 10\,s interval, which matches the default clock skew of 300\,s in ASP.NET's JWT validation [2]; the wording of the 401 also comes from Microsoft's .NET identity library. The second lane of \figref{token}, an open WebSocket that outlived the token, belongs to the live feed (Section~\ref{sec:feed}).

% OWNER: figs/token.tex. Edit only this file.
\fig[wide]{token}{The server accepts an access token for 300\,s after expiry; an open WebSocket outlives it}[Three tests on 24 Nov 2026 aligned at each access token's expiry (x = 0): REST requests with a token issued at 08:16:43 by the server's clock, one WebSocket whose token expired at 07:21:41 KST, and the refresh cookie set at the 05:38:54 login]{%
\begin{tikzpicture}[x=6mm,y=1mm,lab/.style={font=\fontsize{7}{8.4}\selectfont},lane/.style={font=\fontsize{7.5}{9}\selectfont\bfseries,anchor=east,text width=24mm,align=right}]
% ---- lane labels
\node[lane] at (-15.25,0) {REST requests};
\node[lane] at (-15.25,-17) {Open WebSocket};
\node[lane] at (-15.25,-31) {Refresh cookie};
% ---- expiry line
\draw[navy,dashed,line width=0.5pt] (0,-36) -- (0,15.5);
\node[lab,anchor=south east] at (-0.1,12.2) {expiry};
% ---- lane 1: REST
\fill[navy] (-15,-2.5) rectangle (0,2.5);
\fill[navyL] (0,-2.5) rectangle (4.95,2.5);
\fill[pattern=north east lines,pattern color=navyM] (4.95,-2.5) rectangle (5.1167,2.5);
\draw[navyM,line width=0.3pt] (4.95,-2.5) rectangle (5.1167,2.5);
\fill[navyXL] (5.1167,-2.5) rectangle (9,2.5);
\node[lab,text=white] at (-7.5,0) {valid: 900\,s from issue};
\node[lab] at (2.475,0) {accepted};
\node[lab,text=greyD] at (7.06,0) {refused};
\node[lab,anchor=south west,text=greyD] at (-15,2.8) {issued};
% 200 and 401 marks
\fill[navy] (4.95,6) circle (0.75mm); \draw[greyM,dashed,line width=0.3pt] (4.95,5.2) -- (4.95,2.5);
\node[lab,anchor=east] at (4.75,6) {200 at +4:57};
\draw[navy,line width=0.8pt] ($(5.1167,11)+(-0.8mm,-0.8mm)$) -- ($(5.1167,11)+(0.8mm,0.8mm)$) ($(5.1167,11)+(-0.8mm,0.8mm)$) -- ($(5.1167,11)+(0.8mm,-0.8mm)$);
\draw[greyM,dashed,line width=0.3pt] (5.1167,10) -- (5.1167,2.5);
\node[lab,anchor=west] at (5.35,11) {401 at +5:07};
% tolerance bracket (timeout style)
\draw[greyD,line width=0.4pt] (0,-4) -- (0,-5) -- (5,-5) -- (5,-4);
\node[lab,anchor=north,text=greyD] at (2.5,-5.2) {tolerance 300\,s};
% ---- lane 2: WebSocket replies, one tick per reply
\draw[greyM,line width=0.3pt] (-11.42,-17) -- (9,-17);
\foreach \i in {0,...,381}{\draw[navy,line width=0.25pt] ({-11.42+\i*0.05052},-19) -- ({-11.42+\i*0.05052},-15);}
\node[lab,anchor=south west] at (-11.42,-14.6) {1 tick per reply, polled every 3\,s};
\node[lab,anchor=south east] at (7.83,-14.6) {382 replies by the end};
\draw[navy,line width=0.9pt] ($(8.3167,-17)+(-1mm,-1mm)$) -- ($(8.3167,-17)+(1mm,1mm)$) ($(8.3167,-17)+(-1mm,1mm)$) -- ($(8.3167,-17)+(1mm,-1mm)$);
\node[lab,anchor=north east,align=right] at (9.2,-19.6) {new negotiation 07:30 (+8:19),\\ same token: 401};
% ---- lane 3: refresh cookie
\fill[navyM] (-15,-33) rectangle (9,-29);
\draw[navyM,line width=1pt,-{Stealth[length=2mm]}] (9,-31) -- (10.1,-31);
\node[lab,text=white,anchor=west] at (-14.7,-31) {set at login; expires 7 days later, 05:38:54 KST on 1 Dec; renewal leaves this expiry unchanged};
% ---- axis
\draw[greyM,line width=0.4pt] (-15,-38) -- (9,-38);
\foreach \m/\t in {-15/−15,-12/−12,-9/−9,-6/−6,-3/−3,0/0,3/+3,6/+6,9/+9}{\draw[greyM] (\m,-38) -- (\m,-39); \node[lab,text=greyD,anchor=north] at (\m,-39) {\t};}
\node[lab,text=greyD,anchor=north] at (-3,-42.2) {minutes after the access token's expiry};
\end{tikzpicture}
}{Shading: navy where the access token is valid, light navy where it has expired and is still accepted, hatched for the 10-second interval that holds the REST cut-off, pale where REST requests are refused. The tolerance edges come from the second test, made with another token than the WebSocket test, and are placed relative to its expiry; times were measured with the server's \code{Date} header, since our PC clock ran 3\,s ahead. The 401 carried \code{www-authenticate: Bearer error="invalid_token", error_description="The token expired at ..."} (UTC). Each WebSocket tick is one reply on the same connection; the tick positions are spread evenly between the first and last one-per-minute log lines of the test script. The X marks a refused request. Cookie expiry read from \code{Set-Cookie} before and after renewal.}{DART Quant Team}


\lead{Our PC clock ran 3\,s ahead of the server's, so the wrapper times expiry by the server's clock.} In the second test the local clock read 08:36:43 when the server's \code{Date} header read 23:36:40 UTC. A wrapper that compares its own clock with the JWT's \code{exp} inherits that error. The wrapper therefore takes the offset from the \code{Date} header of every reply and uses the median of the last 9 offsets; since the header has 1\,s resolution, it adds 0.5\,s to each sample.

\lead{Renewal happens on demand: the first REST request made 720\,s or more after the token was issued renews it.} There is no timer. Before each REST request the wrapper compares the server's time with the token's \code{iat}; once 720\,s have passed, it calls \code{SilentRefresh} first and then sends the request with the new token. The 180\,s left before the 900\,s expiry cover a slow renewal, and the schedule does not use the 300\,s tolerance, which is a server setting that can change. The live feed needs a token only when it connects or reconnects, because an open WebSocket keeps working after its token expires (Section~\ref{sec:feed}). Concurrent requests share one renewal: in an offline test, 5 simultaneous requests for a token produced 1 \code{SilentRefresh} call.

\lead{A malformed request and a failed login get different errors from different parts of the stack.} A request body with broken JSON (a raw carriage return inside a string) drew an HTML 400 page in nginx's format. Valid JSON with wrong credentials drew a 401 from the API. A 400 points to a defect in the code that built the request and will fail again on retry, so the wrapper raises it as an error without retrying. A 401 is resolved by a renewal or a new login.

\lead{On a 401 the wrapper renews the token once and resends the request; we assume, from ASP.NET Core's documented middleware order, that the server refused the request before acting on it [9].} ASP.NET Core runs its authentication and authorization steps before the handler that serves the request, and a request that fails them receives its 401 without reaching the handler. On that basis a resent request cannot act twice, even when it places an order. We have not observed a 401 on an order request, so this rests on the documentation alone; a second 401 after the renewal is raised as an error.

\lead{Logging in again leaves the existing login session working.} After a second login on the same account, the first login session's cookies still renewed its token, and the browser stayed logged in after a page reload. The wrapper and the web page can therefore run at the same time. Each login, however, adds one more login session that lives 7 days.

\lead{Logging out clears the browser's cookies; the stored cookie jar keeps renewing until its 7-day expiry.} \code{POST /api/Auth/Logout} returned 200 with the body \code{{}} and \code{Set-Cookie} headers that blank \code{refreshToken}, \code{userid} and \code{asset_token} with an expiry date in 1970. We had saved a copy of the cookie jar before the call. After the logout that copy still renewed the token (\code{SilentRefresh} 200), the access token issued before the logout still worked on REST (200), and 2 other login sessions of the same account kept working. The stored jar therefore works as a 7-day credential, whatever happens in the browser, and the wrapper is designed around that.

\lead{The wrapper stores the cookie jar and never the password.} The cookie jar is written to \code{~/.config/}\allowbreak\code{pytimefolio/}\allowbreak\code{session.json} with mode 0600, atomically and with each cookie's expiry, after every login and every renewal. The password is passed to the constructor by the caller, for example from an environment variable, and is used only for a new login. The security consequences of the stored jar are listed in Section~\ref{sec:design}.
\verdict{the wrapper renews at 720\,s by the server's clock, and on a 401 it renews once and resends the request, an order included, since the 401 comes before the order handler (D3, D36).}
\verdict{the wrapper logs in with the caller's password only when renewal fails or the stored cookie jar has less than 1 day left; otherwise it renews from the jar (D5, D17). The jar file keeps each cookie's expiry, so the 1-day rule still applies after a restart. It keeps the jar outside the repository with owner-only permissions and does not call \code{Logout} on shutdown, since the stored jar keeps renewing on the server after the call and the call would blank the local copy, forcing a new login and a new login session at the next start (D16).}
